\documentclass[reqno,11pt]{amsart}

\usepackage{amsmath,amssymb,amsthm,mathtools}
\usepackage[T1]{fontenc}
\usepackage[utf8]{inputenc}
\usepackage{booktabs}
\usepackage[colorlinks=true,linkcolor=blue,citecolor=red,urlcolor=blue]{hyperref}
\usepackage{geometry}
\geometry{a4paper,top=2.5cm,bottom=2.5cm,left=2.8cm,right=2.8cm}

\numberwithin{equation}{section}
\emergencystretch=3em
\theoremstyle{plain}
\newtheorem{theorem}{Theorem}[section]
\newtheorem{proposition}[theorem]{Proposition}
\newtheorem{lemma}[theorem]{Lemma}
\newtheorem{corollary}[theorem]{Corollary}
\theoremstyle{remark}
\newtheorem{remark}[theorem]{Remark}

\newcommand{\eps}{\varepsilon}
\DeclareMathOperator{\tr}{tr}
\DeclareMathOperator{\Real}{Re}
\DeclareMathOperator{\Imag}{Im}

\title[Zeta zeros: unconditional bounds and limits of recent approaches]{Unconditional bounds for zeros of the Riemann zeta function\\ and limits of recent approaches to the Riemann Hypothesis}

\author{Deep Bhattacharjee}
\address{Electro-Gravitational Space Propulsion Laboratory (EGSPL),
Bhubaneswar, Odisha, India.
\mbox{ORCID: \href{https://orcid.org/0000-0003-0466-750X}{0000-0003-0466-750X}}}
\email{itsdeep@live.com}

\author{Priyabrata Mandal}
\address{Department of Mathematics, Bioinformatics and Computer Applications,
Maulana Azad National Institute of Technology Bhopal, Bhopal 462003, India.
\mbox{ORCID: \href{https://orcid.org/0000-0001-6472-6239}{0000-0001-6472-6239}}}
\email{priyabrata@manit.ac.in}
\thanks{Corresponding author: Priyabrata Mandal.}

\date{29 September 2026. Draft, not submitted.}

\begin{document}

\begin{abstract}
This paper does not prove the Riemann Hypothesis. It proves three
unconditional statements (none assumes RH) about the nontrivial zeros of
$\zeta(s)$ and about recent attempts on them.
(i) With the analytic and computer-assisted inputs of Knausg{\aa}rd's recent
argument, a better choice of its free parameters gives
$\liminf_{T\to\infty}N_d(T)/N(T)\ge 0.8369929140\ldots$ for the proportion of
distinct zeros, and the block length we use is the largest the argument
admits.
(ii) For the trace-moment method built on the Gabor-compressed Weil matrix, we
compute exactly the contribution of a pair of zeros off the critical line and
show that any upper bound $\tr\widetilde G^{\,b}\le K\,d\,\ell_1^{\,b}$ with $b$
even forces every extremal, weakly isolated zero $\rho=\beta+i\gamma$ to satisfy
$\beta\le\frac12+\frac1{b\lambda}+o(1)$. For $b=4$ this is a zero-free
statement far beyond what is known. The fourth- and sixth-moment inputs of two
recent preprints claiming $0.7962$ and density one for simple zeros on the
line are of exactly this form. The density-one preprint has been retracted by
its authors for this reason; we show that the $0.7962$ claim depends on the same
tier and should be regarded as unproven.
(iii) We locate the invalid step in X.-J.~Li's claimed proof of RH
(arXiv:0807.0090v10). It is a change of variables that is used in both key
positivity theorems, and a smooth test function shows that the combined
statement those theorems assert is false.
We also recall why a density-one statement, even if it were proved, would not
imply RH.
\end{abstract}

\maketitle

\noindent\textbf{MSC 2020:} 11M26 (primary); 11M06, 11Y35.

\section{Introduction}

Let $N(T)$ count the nontrivial zeros $\rho=\beta+i\gamma$ of $\zeta(s)$ with
$0<\gamma\le T$, with multiplicity, let $N_d(T)$ count the distinct ones, and
let $N_0^{s}(T)$ count those that are simple and lie on the critical line
$\beta=\frac12$. The Riemann Hypothesis (RH) asserts that every nontrivial zero
has $\beta=\frac12$.

\subsection*{What this paper does not prove}
We do not prove RH, and nothing here implies it. We also stress a point that
is easy to overlook. Even the statement $\lim_{T\to\infty}N_0^{s}(T)/N(T)=1$
(``$100\%$ of the zeros are simple and on the line'') would not prove RH. A
proportion statement is insensitive to any set of $o(N(T))$ zeros, and such a
set can still be infinite. Only the statement that \emph{every} nontrivial zero
lies on the line is RH.

\subsection*{What is known}
Selberg \cite{Sel42} proved that a positive proportion of the zeros lie on the
line. Along Levinson's method \cite{Lev74} the refereed records are
Conrey's $2/5$ \cite{Con89} and $5/12$ by Pratt, Robles, Zaharescu and
Zeindler \cite{PRZZ}. Since August 2026 a series of preprints, all
unrefereed, has used Montgomery's pair-correlation method \cite{Mon73} in the
unconditional form of Baluyot, Goldston, Suriajaya and Turnage-Butterbaugh
\cite{BGST24}. They show that more than two thirds of the zeros are simple and on
the line \cite{C26,AF26,Lam26,Wan26}, with computer-assisted refinements to about
$67.3\%$ \cite{TG26,TRM26}. For distinct zeros the best preprint bound not affected by
Section~\ref{s:claims} is $0.83699288\ldots$ \cite{Kna26}.

\subsection*{Results}
Section~\ref{s:distinct} improves the distinct-zero constant of \cite{Kna26}
(Theorem~\ref{thm:distinct}). Section~\ref{s:obstruction} proves an obstruction
for trace-moment methods (Theorem~\ref{thm:obstruction}). Section~\ref{s:claims}
applies it to two recent preprints \cite{YY26a,YY26b} whose headline claims,
$0.7962$ and density one, go far beyond the two-thirds line. Section~\ref{s:li}
examines Li's claimed proof of RH \cite{Li08}. Section~\ref{s:repro} lists the
code that reproduces every computation.

\section{Distinct zeros}\label{s:distinct}

We recall Knausg{\aa}rd's argument \cite{Kna26} only as far as needed. It
combines the unconditional pair-correlation asymptotic \cite{BGST24} with a
mixed-multiplicity matrix inequality and a certified seven-point local
inequality taken from \cite{TRM26}. The imported constants are
\begin{equation}\label{eq:imports}
H=\frac{3362285207}{5000000000},\qquad \delta=\frac{891}{200000},\qquad
\text{spread weight }\frac1{2736}.
\end{equation}
The argument has free parameters $\tau\ge0$, $c$ and $m$. Put
\[
\Delta_m=(m-6)\delta,\quad a=\frac{\Delta_m}{m},\quad
\beta_m=\frac{6(m-6)}{2736\,m},\quad r_h=6c-7-c^2,\quad r_k=4c-2-c^2 .
\]
The proof of \cite[Theorem~1]{Kna26} uses the parameters only through
\begin{align}
&c\ge\max(1+\tau,\,2+\tau/2), \label{eq:A1}\\
&\Delta_m\le\tau^2, \label{eq:A2}\\
&r_h-a\ge0,\qquad r_k-2a\ge0, \label{eq:A3}
\end{align}
and yields
\begin{equation}\label{eq:qm}
\liminf_{T\to\infty}\frac{N_d(T)}{N(T)}\ \ge\ q(m):=\frac{1+H-\beta_m}{2-a}.
\end{equation}
Condition \eqref{eq:A1} is the hypothesis of the mixed-multiplicity Gram
inequality \cite[Thm.~2]{Kna26} and of its combination with the threshold
lemma \cite[Lem.~4]{Kna26}. Condition \eqref{eq:A2} is the hypothesis of the
block dichotomy. Condition \eqref{eq:A3} makes the high-multiplicity and
off-line terms of \cite[(1)--(2)]{Kna26} nonnegative. Knausg{\aa}rd takes
$(\tau,c,m)=(\frac{12}5,\frac{17}5,1298)$ and obtains
$q(1298)=\frac{16260119298029}{19426831050000}=0.83699288145\ldots$.

\begin{theorem}\label{thm:distinct}
With the inputs of \cite{Kna26},
\[
\liminf_{T\to\infty}\frac{N_d(T)}{N(T)}\ \ge\
\frac{62359683640669}{74504434380000}=0.83699291404068\ldots .
\]
\end{theorem}

\begin{proof}
Take $\tau=\frac{2411}{1000}$, $c=\frac{3411}{1000}$ and $m=1310$. Then
$c=1+\tau>2+\tau/2$, which is \eqref{eq:A1}. Next,
$\Delta_m=\frac{145233}{25000}=5.80932<\tau^2=5.812921$, which is
\eqref{eq:A2}. With $a=\frac{145233}{32750000}$ and
$\beta_m=\frac{163}{74670}$ one finds
$r_h-a=\frac{239290417}{131000000}>0$ and $r_k-2a=\frac{5497}{26200000}>0$,
which is \eqref{eq:A3}. The steps of \cite{Kna26} from
\eqref{eq:A1}--\eqref{eq:A3} to \eqref{eq:qm} are stated for general admissible
parameters, and the Lean statements of \cite[\S4]{Kna26} for Theorem~2, Lemma~4
and the counting assembly are parametric. The smoothing error vanishes as
$\eps\to0$ for fixed $(m,\tau)$. Hence \eqref{eq:qm} holds with $m=1310$, and
exact arithmetic gives $q(1310)=\frac{62359683640669}{74504434380000}$.
\end{proof}

The gain over \cite{Kna26} is $3.2588\ldots\cdot10^{-8}$. It is a parameter
refinement, not a new method.

\begin{proposition}\label{prop:limit}
In the argument of \cite{Kna26}, with the same inputs, no admissible
$(\tau,c,m)$ has $m>1310$. If the condition $r_k\ge 2a$ were dropped, the
bound \eqref{eq:qm} would increase to
$q_\infty=(1+H-6/2736)/(2-\delta)=0.8369964390\ldots$ as $m\to\infty$.
\end{proposition}

\begin{proof}
Condition \eqref{eq:A2} forces $\tau\ge\sqrt{\Delta_m}$. If $\tau<2$ then
$\Delta_m<4$ and $m\le903$. For $\tau\ge2$, \eqref{eq:A1} reads $c\ge1+\tau$.
Since $r_k$ decreases for $c>2$, the condition $r_k\ge2a$ is easiest to meet
at $c=1+\sqrt{\Delta_m}$, where it becomes
$1+2\sqrt{\Delta_m}-\Delta_m\ge 2\Delta_m/m$. In this range the left side
decreases and the right side increases with $m$. The inequality holds at
$m=1310$ and fails at $m=1311$ (exact check in \texttt{code/constants.py}).
The condition $r_h\ge a$ is not binding, since $r_h>1$ on $[3,3.5]$. Finally
$q(m)$ increases with $m$, since at leading order the sign of $q'(m)$ is that of
$(1+H-\beta_m)\delta-(2-a)\cdot6/2736>0$, and $q(m)\to q_\infty$.
\end{proof}

\begin{remark}
The binding condition $r_k\ge2a$ comes from a pair of zeros off the line,
which \cite[Lem.~4]{Kna26} charges at cost $c^2$. A negative residue could be
absorbed using a known lower bound $\kappa^{*}$ for the proportion of simple
zeros on the line, since $2k+3h\le(1-\kappa^{*})N$. With $\kappa^{*}>2/5$
\cite{Con89}, or even $\kappa^{*}>0.6731$ \cite{TG26}, a scan over
$c\in[3.4,4.5]$ shows that the penalty always exceeds the gain.
\end{remark}

\section{The off-line pair obstruction for trace-moment methods}\label{s:obstruction}

\subsection{Setting}
We use the notation of \cite[\S2]{YY26a}, which follows \cite{C26}. Let
$l=\log(T/2\pi)$, $\ell_1=l+2\log2-1$, fix $\lambda\in(0,1]$, and put
$L=\lambda l$, $X=e^{L}$ and $d=\lfloor LT/2\pi\rfloor$. Let
$\varrho\in C^3(\mathbb R)$ be nondecreasing with $\varrho=0$ on $(-\infty,0]$
and $\varrho=1$ on $[1,\infty)$. We fix the ramp width $w=1$ (admissible in
\cite{YY26a}, where $1\le w\le L/8$) and put
$\phi(u)=\varrho\bigl((L/2-|u|)/w\bigr)$. For $z\in\mathbb C$ let
\[
\widehat\phi(z)=\int\phi(u)e^{izu}\,du,\qquad
\Phi(z)=\int\phi(u)^2e^{izu}\,du,\qquad
J(y)=\Phi(-iy)=\int\phi(u)^2e^{yu}\,du,
\]
and $a=\frac1L\int\phi^2$. Since $\phi$ is even and real, $J$ is even,
increasing in $|y|$, and $\widehat\phi(\bar z)=\overline{\widehat\phi(z)}$.
With $\tau_k=T+2\pi k/L$ for $0\le k<d$, the compressed Weil matrix is
\[
G=\sum_{\rho}m_\rho\,u_\rho u_\rho^{\mathsf T},\qquad
u_\rho=\bigl(\widehat\phi(\gamma_\rho-\tau_k)\bigr)_{0\le k<d},\qquad
\widetilde G=G/L,
\]
where for $\rho=\frac12+\eta+i\gamma$ we write $\gamma_\rho=\gamma-i\eta$. This
is the zero side of Weil's explicit formula \cite[\S2.1]{YY26a}. Note the
transpose, not the conjugate transpose. The normalised moments are
$\mu_b=\tr\widetilde G^{\,b}/(d\ell_1^{\,b})$. A zero off the line pairs with
$1-\bar\rho$, for which $\gamma_{1-\bar\rho}=\overline{\gamma_\rho}$ and
$u_{1-\bar\rho}=\overline{u_\rho}$.

\begin{lemma}[Gabor--Poisson identity at complex arguments]\label{lem:poisson}
For all $\alpha,\beta\in\mathbb C$,
\[
\sum_{k\in\mathbb Z}\widehat\phi(\alpha-\tau_k)\,\widehat\phi(\beta-\tau_k)
=L\,\Phi(\alpha-\beta).
\]
\end{lemma}

\begin{proof}
For real arguments this is \cite[Lem.~2.2]{C26} (recalled in \cite[\S2.1]{YY26a}).
It follows from Poisson summation, because $\operatorname{supp}\phi$ has length
$L$ and so only the zero frequency survives. Both sides are entire in
$(\alpha,\beta)$. Two integrations by parts give
$|\widehat\phi(x+iy)|\ll_{\varrho} e^{|y|L/2}\min(L,|x|^{-2})$, so the series
converges locally uniformly, and the identity extends by analytic
continuation.
\end{proof}

\begin{lemma}[The block of an off-line pair]\label{lem:block}
Let $\rho_0=\frac12+\eta_0+i\gamma_0$ with $\eta_0>0$ be a zero of
multiplicity $m_0$, and let $u=u_{\rho_0}$. Put $n=u^{*}u=\sum_k|u_k|^2$ and
$s=u^{\mathsf T}u$. The real symmetric matrix
$B_0=m_0(uu^{\mathsf T}+\bar u\bar u^{\mathsf T})$ has rank at most $2$, and its
nonzero eigenvalues are
\[
m_0\Bigl(\Real s\pm\sqrt{n^2-(\Imag s)^2}\Bigr).
\]
If $\gamma_0\in(T+T^{1/2},\,2T-T^{1/2})$, then
$n=L\,J(2\eta_0)+O(LT^{-1})$ and $s=aL^2+O(LT^{-1})$.
\end{lemma}

\begin{proof}
The nonzero eigenvalues of $MM^{\mathsf T}$ with $M=[u\ \bar u]$ are those of
$M^{\mathsf T}M=\begin{pmatrix}s&n\\ n&\bar s\end{pmatrix}$, whose eigenvalues
are $\Real s\pm\sqrt{n^2-(\Imag s)^2}$. By Lemma~\ref{lem:poisson}, the full
sums over $k\in\mathbb Z$ are $L\Phi(z_0-\bar z_0)=L\Phi(-2i\eta_0)=LJ(2\eta_0)$
and $L\Phi(0)=aL^2$, where $z_0=\gamma_0-i\eta_0$. The terms with $k<0$ or
$k\ge d$ have $|\gamma_0-\tau_k|\ge T^{1/2}+2\pi j/L$ with $j\ge0$. By the decay
bound they total $\ll e^{\eta_0L}\sum_{j\ge0}(T^{1/2}+2\pi j/L)^{-4}\ll e^{\eta_0L}L\,T^{-3/2}$.
Since $e^{\eta_0 L}\le X^{1/2}\le T^{1/2}$, this is $O(LT^{-1})$.
\end{proof}

Thus a single off-line pair contributes to $\widetilde G$ a positive eigenvalue
$m_0J(2\eta_0)(1+o(1))$ and a negative one of the same size, with
\begin{equation}\label{eq:Jlower}
J(2\eta_0)\ \ge\ \int_0^{L/2-1}e^{2\eta_0u}\,du
=\frac{e^{2\eta_0(L/2-1)}-1}{2\eta_0}=\frac{X^{\eta_0}e^{-2\eta_0}-1}{2\eta_0}.
\end{equation}
The numerical check in \texttt{code/offline\_pair.py} confirms
Lemma~\ref{lem:poisson} at complex arguments to $2\cdot10^{-9}$ and
Lemma~\ref{lem:block} to ten digits.

\subsection{The obstruction}
Let $I'=(T-T^{1/2},\,2T+T^{1/2}]$. For an off-line zero $\rho_0$ with
$\gamma_0\in I'$ define its \emph{interference}
\[
\mathcal I(\rho_0)=\sum_{\rho}{}^{'}\,m_\rho\min\bigl(1,|\gamma_\rho-\gamma_0|^{-4}\bigr),
\]
where $\sum'$ runs over the off-line zeros $\rho\notin\{\rho_0,1-\bar\rho_0\}$
with $\gamma_\rho\in I'$. Call $\rho_0$ \emph{extremal} if
$\beta_\rho\le\beta_0$ for every zero $\rho$ with $\gamma_\rho\in I'$.

\begin{theorem}\label{thm:obstruction}
There is a constant $c_\varrho>0$ depending only on $\varrho$ with the
following property. Let $T$ be large and let $\rho_0=\frac12+\eta_0+i\gamma_0$
be an extremal zero with $\eta_0>0$,
$\gamma_0\in(T+T^{1/2},2T-T^{1/2})$ and $\mathcal I(\rho_0)\le c_\varrho$. Then
\[
\|\widetilde G\|_{\mathrm{op}}\ \ge\ \tfrac12J(2\eta_0)-3L-1,
\qquad\text{hence}\qquad
\tr\widetilde G^{\,b}\ \ge\ \bigl(\tfrac12J(2\eta_0)-3L-1\bigr)^{b}
\ \text{ for every even } b.
\]
\end{theorem}

\begin{proof}
Split $G=B_0+P+Q+E$, where $B_0$ is the block of $\{\rho_0,1-\bar\rho_0\}$ (the
two zeros have the same multiplicity), $P$ collects the zeros on the line with
$\gamma\in I'$, $Q$ collects the other off-line zeros with $\gamma\in I'$
(grouped in pairs), and $E$ collects the zeros with $\gamma\notin I'$. The matrix
$P$ is positive semidefinite, because $u_\rho$ is real for $\beta_\rho=\frac12$.
By the tail estimate of \cite[\S2.2]{YY26a} (that is, \cite[Prop.~4.2]{C26}),
$\|E\|_{\mathrm{op}}/L\le\theta_0\ll lT^{\lambda/2-1}\le1$.

Let $x$ be a unit eigenvector of $B_0$ for its largest eigenvalue. By
Lemma~\ref{lem:block}, $x^{\mathsf T}B_0x\ge m_0(n-2|s|)\ge LJ(2\eta_0)-2L^2-1$.
The vector $x$ is real and lies in the span of $u$ and $\bar u$. Write
$x=\alpha u+\beta\bar u$. Then
$1=|x|^2\ge(|\alpha|^2+|\beta|^2)(n-|s|)$, because the Hermitian product of $u$
and $\bar u$ equals $\bar s$.

If $J(2\eta_0)\le6L+2$ the right side of the claim is at most $0$ and
there is nothing to prove, so assume $J(2\eta_0)>6L+2$. Since $a\le1$, this
gives $n-|s|\ge LJ(2\eta_0)-L^2-1\ge\frac45LJ(2\eta_0)$ for large $T$.

Now let $\rho$ be one of the other off-line zeros, with block
$B_\rho=m_\rho(u_\rho u_\rho^{\mathsf T}+\bar u_\rho\bar u_\rho^{\mathsf T})$. Then
$|x^{\mathsf T}B_\rho x|\le2m_\rho|u_\rho^{\mathsf T}x|^2$, and
\[
|u_\rho^{\mathsf T}x|^2\le\frac{2\bigl(|u_\rho^{\mathsf T}u|+|u_\rho^{\mathsf T}\bar u|\bigr)^2}{n-|s|}.
\]
We compare the inner products with the full sums over $k\in\mathbb Z$, which
Lemma~\ref{lem:poisson} evaluates as $L\Phi(\gamma_\rho-z_0)$ and
$L\Phi(\gamma_\rho-\bar z_0)$. Every omitted term has
$|\gamma_0-\tau_k|\ge T^{1/2}$. Using the decay bound for both factors, the
omitted terms total $t_\rho\ll e^{\eta_0L}L^{3/2}T^{-1}$. The imaginary parts
of $\gamma_\rho-z_0$ and $\gamma_\rho-\bar z_0$ have absolute value at most
$2\eta_0<1$, by extremality. Integrating by parts twice gives, for real $x$ and
$|y|\le1$,
\[
|\Phi(x+iy)|\le\frac1{x^2}\Bigl(y^2J(|y|)+(2|y|\,\|(\phi^2)'\|_\infty+\|(\phi^2)''\|_\infty)\int_{\text{ramps}}e^{|y||u|}du\Bigr)
\le\frac{C_\varrho\,(J(|y|)+1)}{x^2}.
\]
Here the ramps are the two unit intervals where $\phi<1$, and we used
$\int_{\text{ramps}}e^{|y||u|}du\le2e^{|y|L/2}\le2e^{|y|}(|y|J(|y|)+1)$, which
follows from \eqref{eq:Jlower}. Together with the trivial bound
$|\Phi(x+iy)|\le J(|y|)$ and the monotonicity of $J$,
\[
|u_\rho^{\mathsf T}u|+|u_\rho^{\mathsf T}\bar u|\le
2C'_\varrho L\bigl(J(2\eta_0)+1\bigr)\min(1,\delta_\rho^{-2})+2t_\rho,\qquad
\delta_\rho=|\gamma_\rho-\gamma_0|,
\]
with $C'_\varrho=\max(1,C_\varrho)$. Using $(p+q)^2\le2p^2+2q^2$,
$(J+1)^2\le1.1J^2$ and $n-|s|\ge\frac45LJ$, and summing over the pairs in $Q$,
\[
|x^{\mathsf T}Qx|\le 44\,C_\varrho'^2\,L\,J(2\eta_0)\,\mathcal I(\rho_0)
+O\Bigl(\sum_{\rho}{}^{'}\frac{m_\rho t_\rho^2}{LJ(2\eta_0)}\Bigr).
\]
There are $\ll T\log T$ zeros with $\gamma\in I'$, counted with
multiplicity. Since $J(2\eta_0)\gg e^{\eta_0L}$ and $e^{\eta_0L}\le T^{1/2}$,
the error term is $\ll L^2(\log T)e^{\eta_0 L}T^{-1}\le1$.

Choose $c_\varrho=1/(160\,C_\varrho'^2)$, so that
$|x^{\mathsf T}Qx|\le0.3\,LJ(2\eta_0)+1$. Then
\[
x^{\mathsf T}Gx\ \ge\ x^{\mathsf T}B_0x-|x^{\mathsf T}Qx|-\|E\|_{\mathrm{op}}
\ \ge\ 0.7\,LJ(2\eta_0)-2L^2-L-2\ \ge\ \tfrac12LJ(2\eta_0)-3L^2-L,
\]
because $x^{\mathsf T}Px\ge0$. Dividing by $L$ gives the bound for
$\|\widetilde G\|_{\mathrm{op}}\ge x^{\mathsf T}\widetilde Gx$. Finally
$\tr\widetilde G^{\,b}=\sum_i\lambda_i^b\ge\|\widetilde G\|_{\mathrm{op}}^b$
for even $b$.
\end{proof}

\begin{corollary}\label{cor:zerofree}
Fix $\lambda\in(0,1]$, an even $b\ge2$ and $K>0$, and suppose that
$\tr\widetilde G^{\,b}\le K\,d\,\ell_1^{\,b}$ for all large $T$. Then every
extremal zero $\rho_0$ as in Theorem~\ref{thm:obstruction} with
$\mathcal I(\rho_0)\le c_\varrho$ satisfies
\[
\beta_0\ \le\ \frac12+\frac1{b\lambda}+O\Bigl(\frac{\log\log T}{\log T}\Bigr).
\]
\end{corollary}

\begin{proof}
By Theorem~\ref{thm:obstruction} and \eqref{eq:Jlower},
$(X^{\eta_0}e^{-2\eta_0}-1)/(4\eta_0)\le(Kd)^{1/b}\ell_1+3L+1$. Since
$d\le LT/2\pi$, $\ell_1\ll\log T$ and $X=(T/2\pi)^{\lambda}$, taking logarithms
gives $\lambda\eta_0\log T\le\frac1b\log T+O(\log\log T)$.
\end{proof}

\begin{remark}[What the corollary means]\label{rem:meaning}
For $b=2$ the conclusion is $\beta_0\le\frac12+\frac1{2\lambda}+o(1)$, which is
empty because $\beta_0<1$ and $\lambda\le1$. This is consistent with the second
moment $\mu_2\to\frac43$ being a theorem \cite{C26,BGST24}: Montgomery's
bandwidth $\lambda\le1$ is exactly what keeps it clear of the obstruction.
For $b=4$ and $\lambda$ close to $1$ the conclusion is
$\beta_0\le\frac34+\eps$. Nothing of this kind is known: the best
unconditional zero-free region is of Vinogradov--Korobov type,
$\beta<1-c(\log\gamma)^{-2/3}(\log\log\gamma)^{-1/3}$ \cite{For02}. For $b=6$
it would be $\beta_0\le\frac23+\eps$. Odd moments are not covered, because the
two eigenvalues of Lemma~\ref{lem:block} cancel to leading order in
$\tr\widetilde G^{\,b}$ for odd $b$. This agrees with the retraction notice
of \cite{YY26b}, which speaks of higher \emph{even} moments.
\end{remark}

\begin{remark}[Limits of Theorem~\ref{thm:obstruction}]
The interference hypothesis excludes clusters of off-line zeros near
$\gamma_0$. Such a cluster could in principle produce cancellation in
$x^{\mathsf T}Qx$, and we do not rule this out. The authors of \cite{YY26b}
state that a reviewer's argument (``Proposition R''), which they verified,
shows in general that finite unconditional bounds on higher even moments would
force zero-free regions of power-law strength. That argument is not public,
and we do not rely on it. Theorem~\ref{thm:obstruction} proves the
mechanism rigorously in the isolated case, which suffices for our conclusions
below.
\end{remark}

\section{Consequences for two recent claims}\label{s:claims}

\subsection{The density-one preprint}
The preprint \cite{YY26b} (version~0.92, August 2026) claimed unconditionally that
$\lim N_0^{s}(T)/N(T)=1$. It evaluated the trace moments of $\widetilde G$
through order fourteen and consumed them in a Christoffel-function count. On
3~September 2026 its authors retracted the central claims (repository commit
\texttt{ac85152}, file \texttt{RETRACTION.md}), stating that ``finite
unconditional bounds on higher even moments would already force zero-free
regions of power-law strength'' and that ``the gap was located; the claims do
not survive''. Corollary~\ref{cor:zerofree} gives an independent, rigorous
account of the mechanism in the isolated case.

\subsection{\texorpdfstring{The $0.7962$ preprint}{The 0.7962 preprint}}
The preprint \cite{YY26a} (repository commit \texttt{d85bddf}, 17 August 2026;
Zenodo DOI 10.5281/zenodo.21975236) claims $\liminf N_0^{s}/N\ge0.7962$ and
$\liminf N_d/N\ge0.8981$. It grades these as ``certified candidates''. It is
not marked as retracted. Its rungs and their inputs are as follows (line
numbers refer to \texttt{paper.tex} in that repository).
\begin{itemize}
\item The rung $0.6916$ uses the upper bound $\limsup R(1)\le0.0066$, that is
$\mu_4\le\frac{13}4+0.0399+o(1)$ (the one-sided fourth-moment theorem,
lines~882--897).
\item The rung $13/18$ uses ``$m_4\to13/4$ two-sidedly'' (Lemma~D in truncated
form, lines~1365--1376, and its corollary, line~1427).
\item The headline $0.7962$ additionally uses $m_6\le M_6=\frac{12809}{1260}$
(section ``Consumption and the proof of the main theorem'', lines~1823--1835).
\end{itemize}
The paper notes that the count is monotone in each even moment (line~544), so
every rung above the pair-correlation ceiling $0.68185$ of \cite{C26}, which
the paper recalls at line~146, consumes an \emph{upper} bound on $\mu_4$ and,
for the headline, on $\mu_6$. By Corollary~\ref{cor:zerofree} with $b=4$
(respectively $b=6$), each of these inputs, if proved unconditionally, would
imply that $\zeta$ has no extremal, weakly isolated zeros with $\beta>\frac34+\eps$
(respectively $\beta>\frac23+\eps$) for large height. That is far beyond current
knowledge. The paper's own objects ($\widetilde G$, the moment $\mu_4\to13/4$)
are literally those of the retracted tier of \cite{YY26b}; compare line~533 with
the retraction notice. The paper states that ``proportion statements are
insensitive to $o(N)$ off-line zeros'' (line~191). Lemma~\ref{lem:block} shows
that the trace moments of order $b\ge4$ are not: a single off-line pair changes
$\tr\widetilde G^{\,4}$ by $\asymp X^{4\eta_0}/\eta_0^4$.

We conclude that the claims $0.6916$, $13/18$, $0.7962$ and the corresponding
distinct-zero claims of \cite{YY26a} should be regarded as \emph{unproven}.
We have not located the specific line where the proof of Lemma~D fails. What
Corollary~\ref{cor:zerofree} shows is that, unless that proof contains a gap,
it proves a zero-free statement of a strength no known method reaches.

\section{Li's claimed proof of the Riemann Hypothesis}\label{s:li}

X.-J.~Li's arXiv preprint \cite{Li08} (first version July 2008, version~10
October 2025, listed under math.GM) claims to prove RH. It works in Connes'
semilocal adelic framework \cite{Con99}. The chain is as follows.
Theorem~1.1 (cited) gives $\tr(T_h)=\Delta(h)-\hat h(0)-\hat h(1)$ with
$\Delta(h)=\sum_\rho\hat h(\rho)$ and $\hat h(s)=\hat g(s)\hat g(1-s)$.
Theorem~1.2 relates special test functions $g_{n,\eps}$ to the Li coefficients
$\lambda_n$. Theorem~1.3 asserts that the trace of $T_h$ on
$E_S(Q_\Lambda^{\perp})$ vanishes, and Theorem~1.4 that the trace on
$E_S(Q_\Lambda)$ is nonnegative. Li's criterion \cite{Li97} then gives RH.

\subsection{The invalid step}
Section~4 of \cite{Li08} ends with an absolutely convergent sum over the
infinite unit group $O_S^{*}$, equation~(4.13). The paper substitutes
$x\mapsto\xi^{-1}x$ and then $u\mapsto u\xi$, and asserts that every term of
(4.14) is the same number. Since the sum converges, that number must be $0$,
which is (4.15). The substitution in $\int_{C_S}\cdots d^{\times}x$ is valid only
if the integrand is a function on $C_S=J_S/O_S^{*}$, that is, invariant under
multiplication of $x$ by units. The integrand contains $\Psi_S(x\xi v)$, which
is not invariant. For $S=\{\infty,2\}$ the element $\frac12$ lies in $O_S^{*}$,
and replacing $v$ by $v/2$ changes the normalised local factor
$\int_{\mathbb Z_2^{*}}\psi_2(-xv)\,d^{\times}x$ from $+1$ to $-1$
(\texttt{audit/code/padic\_check.py}, which also checks Lemma~2.2 of
\cite{Li08} for $p=2,3,5,7$). So the terms of (4.14) are not equal. Section~5
makes the same move at (5.8)--(5.10), so Theorems~1.3 and~1.4 are both
unproven.

\subsection{The argument proves too much}
The proofs in Sections~4 and~5 of \cite{Li08} use only that $g$ is real,
smooth and supported in $[1-\eps,\mu_\eps]$. If Theorems~1.1, 1.3 and~1.4 held
together for all such $g$, they would give
\begin{equation}\label{eq:li-star}
\sum_\rho\hat g(\rho)\,\hat g(1-\rho)\ \ge\ \hat h(0)+\hat h(1)=2\,\hat g(0)\,\hat g(1).
\end{equation}

\begin{proposition}[numerical]\label{prop:li}
Let $g(x)=G(\log x)$, where $G(u)=\exp\bigl(-1/(1-y^2)\bigr)$ with
$y=2u/3-1$ on $[0,3]$. Then $2\hat g(0)\hat g(1)=4.7287408\ldots$, while
$\sum_\rho\hat g(\rho)\hat g(1-\rho)=0.000182557\ldots$. So
\eqref{eq:li-star} fails by a factor of about $2.6\cdot10^4$.
\end{proposition}

The zero sum is computed in two independent ways. The first uses the first
$50$ zeros, with a tail bound $\le1.3\cdot10^{-7}$ that holds whether or not
the remaining zeros are on the line. The second uses Weil's explicit formula
in the form of \cite{Bom00}, which involves primes and the archimedean term
only. The two agree to $1.7\cdot10^{-12}$
(\texttt{audit/code/weil\_counterexample.py}). Since the proofs of
Theorems~1.3 and~1.4 do not distinguish this $g$ from the functions $g_{n,\eps}$,
they are not valid for $g_{n,\eps}$ either. When $\hat h(0)=\hat h(1)=0$, the
inequality $\tr(T_h)\ge0$ is Weil's positivity criterion, which is equivalent
to RH. The paper therefore does not establish RH.

\begin{remark}
The Li coefficients $\lambda_1,\dots,\lambda_{30}$, computed from derivatives
of $\log\xi$ at $s=1$ without using zeros, are all positive
(\texttt{audit/code/li\_coefficients.py}). This is consistent with RH and
is known far beyond $n=30$. No finite range of $n$ implies RH.
\end{remark}

\subsection{A related claim}
The 2018 manuscript \cite{Esw18} reduces RH to Littlewood's criterion
$L(N)=\sum_{n\le N}\lambda(n)\ll N^{1/2+\eps}$ for the Liouville function. It
then argues that $\lambda(n)$ behaves like a sequence of fair coin tosses. The
properties it establishes (equal densities of $\pm1$, non-periodicity,
finite-range statistics) are compatible with partial sums as large as
$N^{0.99}$, and equal density is equivalent to the prime number theorem, which
gives only $L(N)=o(N)$. The required bound is equivalent to RH and is in effect
assumed, so this is not a proof.

\section{Reproducibility}\label{s:repro}

All code is in the accompanying repository and runs with Python~3, numpy,
mpmath and sympy.
\begin{itemize}
\item \texttt{code/constants.py}: exact rational verification of
Theorem~\ref{thm:distinct} and Proposition~\ref{prop:limit}.
\item \texttt{code/offline\_pair.py}: Lemmas~\ref{lem:poisson}
and~\ref{lem:block} at complex arguments, and the thresholds of
Corollary~\ref{cor:zerofree} at finite height (output in
\texttt{code/offline\_pair.txt}).
\item \texttt{audit/code/}: \texttt{weil\_counterexample.py}
(Proposition~\ref{prop:li}), \texttt{padic\_check.py},
\texttt{li\_coefficients.py} and the zero cache.
\end{itemize}

\subsection*{Status of the inputs}
Theorem~\ref{thm:distinct} depends on the unrefereed certificate of
\cite{TRM26} and on \cite{BGST24}, exactly as \cite{Kna26} does.
Theorem~\ref{thm:obstruction} depends only on the explicit construction of
\cite{C26,YY26a} and on classical estimates. Proposition~\ref{prop:li} is a
numerical computation with an explicit tail bound.

\subsection*{Use of AI tools}
The analyses, the code and a draft of this text were produced with the
assistance of an AI system (Anthropic Claude) under the authors' direction.
The authors are responsible for the content.

\begin{thebibliography}{99}

\bibitem{AF26} L.~Alp\"oge and R.~Furman, \emph{More than two thirds of the
zeta zeros are simple and on the critical line}, arXiv:2608.13637, 2026.

\bibitem{BGST24} S.~A.~C.~Baluyot, D.~A.~Goldston, A.~I.~Suriajaya and
C.~L.~Turnage-Butterbaugh, \emph{An unconditional Montgomery theorem for pair
correlation of zeros of the Riemann zeta-function}, Acta Arith.\ \textbf{214}
(2024), 357--376.

\bibitem{Bom00} E.~Bombieri, \emph{Remarks on Weil's quadratic functional in
the theory of prime numbers, I}, Rend.\ Mat.\ Acc.\ Lincei (9) \textbf{11}
(2000), 183--233.

\bibitem{C26} Claude (Anthropic), \emph{More than two thirds of the zeros of
the Riemann zeta function lie on the critical line}, preprint, 2026.

\bibitem{Con89} J.~B.~Conrey, \emph{More than two fifths of the zeros of the
Riemann zeta function are on the critical line}, J.\ Reine Angew.\ Math.\
\textbf{399} (1989), 1--26.

\bibitem{Con99} A.~Connes, \emph{Trace formula in noncommutative geometry and
the zeros of the Riemann zeta function}, Selecta Math.\ (N.S.) \textbf{5}
(1999), 29--106.

\bibitem{Esw18} K.~Eswaran, \emph{The final and exhaustive proof of the
Riemann Hypothesis from first principles}, manuscript (ResearchGate), May 2018.

\bibitem{For02} K.~Ford, \emph{Vinogradov's integral and bounds for the
Riemann zeta function}, Proc.\ London Math.\ Soc.\ (3) \textbf{85} (2002),
565--633.

\bibitem{Kna26} K.~M.~Knausg{\aa}rd, \emph{More than 83.69\% of the zeros of
the Riemann zeta function are distinct}, arXiv:2609.33043, 2026.

\bibitem{Lam26} Y.~Lamzouri, \emph{A new proof that more than $2/3$ of the
zeros of the Riemann zeta function are simple and on the critical line},
arXiv:2609.02882, 2026.

\bibitem{Lev74} N.~Levinson, \emph{More than one third of zeros of Riemann's
zeta-function are on $\sigma=1/2$}, Adv.\ Math.\ \textbf{13} (1974), 383--436.

\bibitem{Li97} X.-J.~Li, \emph{The positivity of a sequence of numbers and the
Riemann hypothesis}, J.\ Number Theory \textbf{65} (1997), 325--333.

\bibitem{Li08} X.-J.~Li, \emph{A proof of the Riemann hypothesis},
arXiv:0807.0090v10, 2025.

\bibitem{Mon73} H.~L.~Montgomery, \emph{The pair correlation of zeros of the
zeta function}, Proc.\ Sympos.\ Pure Math.\ \textbf{24}, AMS, 1973, 181--193.

\bibitem{PRZZ} K.~Pratt, N.~Robles, A.~Zaharescu and D.~Zeindler, \emph{More
than five-twelfths of the zeros of $\zeta$ are on the critical line}, Res.\
Math.\ Sci.\ \textbf{7} (2020), Paper No.~2.

\bibitem{Sel42} A.~Selberg, \emph{On the zeros of Riemann's zeta-function},
Skr.\ Norske Vid.-Akad.\ Oslo I (1942), no.~10.

\bibitem{TG26} tawanerguo-cn/zeta-simple-zeros, \emph{A certified
unconditional 67.3192911473\% lower bound for simple zeros of the Riemann
zeta function}, research repository, 2026,
\url{https://github.com/tawanerguo-cn/zeta-simple-zeros}.

\bibitem{TRM26} trmdy/zeta-simple-zeros-673137, commit 1610b97, 2026,
\url{https://github.com/trmdy/zeta-simple-zeros-673137}.

\bibitem{Wan26} B.~Wang, \emph{Proportions of the non-trivial zeros of the
Riemann zeta function}, arXiv:2609.24167, 2026.

\bibitem{YY26a} H.~Yang and S.~Yang, \emph{More than 0.7962 of the zeros of
the Riemann zeta function are simple and on the critical line}, preprint,
2026, Zenodo DOI 10.5281/zenodo.21975236;
\url{https://github.com/JoshuaHKU/zeta-0.7947-reproduction}, commit d85bddf.

\bibitem{YY26b} H.~Yang and S.~Yang, \emph{Almost all zeros of the Riemann
zeta function are simple and on the critical line}, preprint v0.92, 2026,
Zenodo DOI 10.5281/zenodo.22065921; retracted 3 September 2026,
\url{https://github.com/JoshuaHKU/zeta-density-one-reproduction}.

\end{thebibliography}

\end{document}
